% Options for packages loaded elsewhere
\PassOptionsToPackage{unicode}{hyperref}
\PassOptionsToPackage{hyphens}{url}
%
\documentclass[
  11pt,
]{article}
\usepackage{amsmath,amssymb}
\usepackage{iftex}
\ifPDFTeX
  \usepackage[T1]{fontenc}
  \usepackage[utf8]{inputenc}
  \usepackage{textcomp} % provide euro and other symbols
\else % if luatex or xetex
  \usepackage{unicode-math} % this also loads fontspec
  \defaultfontfeatures{Scale=MatchLowercase}
  \defaultfontfeatures[\rmfamily]{Ligatures=TeX,Scale=1}
\fi
\usepackage{lmodern}
\ifPDFTeX\else
  % xetex/luatex font selection
\fi
% Use upquote if available, for straight quotes in verbatim environments
\IfFileExists{upquote.sty}{\usepackage{upquote}}{}
\IfFileExists{microtype.sty}{% use microtype if available
  \usepackage[]{microtype}
  \UseMicrotypeSet[protrusion]{basicmath} % disable protrusion for tt fonts
}{}
\makeatletter
\@ifundefined{KOMAClassName}{% if non-KOMA class
  \IfFileExists{parskip.sty}{%
    \usepackage{parskip}
  }{% else
    \setlength{\parindent}{0pt}
    \setlength{\parskip}{6pt plus 2pt minus 1pt}}
}{% if KOMA class
  \KOMAoptions{parskip=half}}
\makeatother
\usepackage{xcolor}
\usepackage[margin=1in]{geometry}
\usepackage{color}
\usepackage{fancyvrb}
\newcommand{\VerbBar}{|}
\newcommand{\VERB}{\Verb[commandchars=\\\{\}]}
\DefineVerbatimEnvironment{Highlighting}{Verbatim}{commandchars=\\\{\}}
% Add ',fontsize=\small' for more characters per line
\newenvironment{Shaded}{}{}
\newcommand{\AlertTok}[1]{\textcolor[rgb]{1.00,0.00,0.00}{\textbf{#1}}}
\newcommand{\AnnotationTok}[1]{\textcolor[rgb]{0.38,0.63,0.69}{\textbf{\textit{#1}}}}
\newcommand{\AttributeTok}[1]{\textcolor[rgb]{0.49,0.56,0.16}{#1}}
\newcommand{\BaseNTok}[1]{\textcolor[rgb]{0.25,0.63,0.44}{#1}}
\newcommand{\BuiltInTok}[1]{\textcolor[rgb]{0.00,0.50,0.00}{#1}}
\newcommand{\CharTok}[1]{\textcolor[rgb]{0.25,0.44,0.63}{#1}}
\newcommand{\CommentTok}[1]{\textcolor[rgb]{0.38,0.63,0.69}{\textit{#1}}}
\newcommand{\CommentVarTok}[1]{\textcolor[rgb]{0.38,0.63,0.69}{\textbf{\textit{#1}}}}
\newcommand{\ConstantTok}[1]{\textcolor[rgb]{0.53,0.00,0.00}{#1}}
\newcommand{\ControlFlowTok}[1]{\textcolor[rgb]{0.00,0.44,0.13}{\textbf{#1}}}
\newcommand{\DataTypeTok}[1]{\textcolor[rgb]{0.56,0.13,0.00}{#1}}
\newcommand{\DecValTok}[1]{\textcolor[rgb]{0.25,0.63,0.44}{#1}}
\newcommand{\DocumentationTok}[1]{\textcolor[rgb]{0.73,0.13,0.13}{\textit{#1}}}
\newcommand{\ErrorTok}[1]{\textcolor[rgb]{1.00,0.00,0.00}{\textbf{#1}}}
\newcommand{\ExtensionTok}[1]{#1}
\newcommand{\FloatTok}[1]{\textcolor[rgb]{0.25,0.63,0.44}{#1}}
\newcommand{\FunctionTok}[1]{\textcolor[rgb]{0.02,0.16,0.49}{#1}}
\newcommand{\ImportTok}[1]{\textcolor[rgb]{0.00,0.50,0.00}{\textbf{#1}}}
\newcommand{\InformationTok}[1]{\textcolor[rgb]{0.38,0.63,0.69}{\textbf{\textit{#1}}}}
\newcommand{\KeywordTok}[1]{\textcolor[rgb]{0.00,0.44,0.13}{\textbf{#1}}}
\newcommand{\NormalTok}[1]{#1}
\newcommand{\OperatorTok}[1]{\textcolor[rgb]{0.40,0.40,0.40}{#1}}
\newcommand{\OtherTok}[1]{\textcolor[rgb]{0.00,0.44,0.13}{#1}}
\newcommand{\PreprocessorTok}[1]{\textcolor[rgb]{0.74,0.48,0.00}{#1}}
\newcommand{\RegionMarkerTok}[1]{#1}
\newcommand{\SpecialCharTok}[1]{\textcolor[rgb]{0.25,0.44,0.63}{#1}}
\newcommand{\SpecialStringTok}[1]{\textcolor[rgb]{0.73,0.40,0.53}{#1}}
\newcommand{\StringTok}[1]{\textcolor[rgb]{0.25,0.44,0.63}{#1}}
\newcommand{\VariableTok}[1]{\textcolor[rgb]{0.10,0.09,0.49}{#1}}
\newcommand{\VerbatimStringTok}[1]{\textcolor[rgb]{0.25,0.44,0.63}{#1}}
\newcommand{\WarningTok}[1]{\textcolor[rgb]{0.38,0.63,0.69}{\textbf{\textit{#1}}}}
\setlength{\emergencystretch}{3em} % prevent overfull lines
\providecommand{\tightlist}{%
  \setlength{\itemsep}{0pt}\setlength{\parskip}{0pt}}
\setcounter{secnumdepth}{-\maxdimen} % remove section numbering
\pdfinfoomitdate=1
\usepackage[T1]{fontenc}
\usepackage[utf8]{inputenc}
\usepackage{microtype}
\usepackage{amsmath,amssymb,mathtools}
\usepackage{booktabs}
\usepackage{geometry}
\usepackage{hyperref}
\usepackage{enumitem}
\usepackage{fancyhdr}
\usepackage{titlesec}
\geometry{margin=1in}
\setlength{\parindent}{0pt}
\setlength{\parskip}{0.55em}
\setlist{itemsep=0.25em, topsep=0.3em}
\hypersetup{hidelinks,breaklinks=true,pdfcreator={LaTeX}}
\pagestyle{fancy}
\fancyhf{}
\fancyhead[L]{Industrial Anomaly Detection}
\fancyhead[R]{VisA PCB1 Technical Analysis}
\fancyfoot[C]{\thepage}
\setlength{\headheight}{14pt}
\titleformat{\section}{\Large\bfseries}{\thesection}{0.75em}{}
\titleformat{\subsection}{\large\bfseries}{\thesubsection}{0.75em}{}
\titleformat{\subsubsection}{\normalsize\bfseries}{\thesubsubsection}{0.75em}{}
\ifLuaTeX
  \usepackage{selnolig}  % disable illegal ligatures
\fi
\usepackage{bookmark}
\IfFileExists{xurl.sty}{\usepackage{xurl}}{} % add URL line breaks if available
\urlstyle{same}
\hypersetup{
  pdftitle={Trustworthy Industrial Anomaly Detection Under Distribution Shift},
  hidelinks,
  pdfcreator={LaTeX via pandoc}}

\title{Trustworthy Industrial Anomaly Detection Under Distribution
Shift}
\usepackage{etoolbox}
\makeatletter
\providecommand{\subtitle}[1]{% add subtitle to \maketitle
  \apptocmd{\@title}{\par {\large #1 \par}}{}{}
}
\makeatother
\subtitle{Technical analysis of the VisA PCB1 implementation}
\author{}
\date{}

\begin{document}
\maketitle

\subsection{Abstract}\label{abstract}

This report provides a technical analysis of the provided Python
implementation for trustworthy industrial anomaly detection on the VisA
PCB1 category. The source is a Colab-oriented, cell-structured Python
script focused on leakage-free anomaly scoring, calibration, robustness
stress-testing, and multi-scale representation. The implementation
reconstructs the PCB1 dataset from verified Parquet files, extracts frozen image representations with
pretrained ResNet-50 and DINOv2 ViT-S/14 backbones, constructs a normal
reference bank from training data, assigns anomaly scores using
cosine-distance k-nearest-neighbor scoring, and maps those scores to
probabilities using logistic and isotonic calibration learned without
access to the real test defects.

The analysis follows the code as implemented rather than replacing it
with an idealized textbook pipeline. It focuses on the separation of
reference and calibration data, the use of synthetic calibration
corruptions, the mathematical form of the anomaly score and calibration
functions, the adaptive expected calibration error implementation,
risk-coverage analysis, hyperparameter sensitivity, robustness evaluation
under image corruption, and the packaging of generated artifacts. The analysis also identifies implementation
constraints that matter for scientific interpretation, including partial
rather than strict determinism, environment assumptions, feature-cache
invalidation risks, the mismatch between the \texttt{ece15} output name
and the ten-bin implementation, and the altered class prevalence induced
by the synthetic calibration construction.

No experimental results, figures, screenshots, or plots are reproduced
in this report. The document is intentionally focused on computational
understanding, methodology, mathematical formulation, software
structure, and reproducibility.

\section{1. Introduction}\label{introduction}

Industrial visual anomaly detection addresses a setting in which a model
is expected to recognize defective or otherwise unusual observations
even though examples of every possible defect are not available during
model development. The provided implementation approaches this problem
as a one-class reference-learning task: normal training images provide a
representation of expected appearance, and test images are assigned
anomaly scores according to their feature-space distance from that
normal reference set.

The implementation is explicitly framed as a trustworthy
anomaly-detection pipeline under distribution shift. That framing is
reflected in the code structure. The script does not simply train a
classifier on the available labels. Instead, it separates three roles
that are often conflated in anomaly-detection experiments: a normal
reference bank for defining the nominal appearance manifold, a separate
calibration pool for learning a score-to-probability mapping, and a
held-out test split containing real industrial defects for final
evaluation. Synthetic defects are introduced only in the calibration
pool so that probability calibration can be learned without using the
real test anomalies.

The computational pipeline combines standard computer-vision backbones
with non-parametric anomaly scoring. ResNet-50 is used as a frozen
feature extractor after removal of its classification layer, while
DINOv2 ViT-S/14 is used through its feature-extraction interface. For
DINOv2, the implementation concatenates the normalized class token with
the mean of the normalized patch tokens, producing a representation that
combines global and aggregated local information. All resulting
embeddings are L2-normalized, making the subsequent inner product
equivalent to cosine similarity. The k-nearest-neighbor anomaly score is
then defined as the average cosine distance to the nearest elements of
the normal reference bank.

Calibration is treated as a separate modeling layer. The raw anomaly
score is not itself interpreted as a probability. Instead, the
implementation fits both a parametric logistic mapping and a
non-parametric isotonic mapping on a calibration set consisting of clean
normal images and several synthetic defect variants. This distinction is
scientifically useful because discrimination and calibration answer
different questions: discrimination asks whether anomalies receive
larger scores than normal observations, whereas calibration asks whether
numerical probabilities have the intended probabilistic interpretation.

The rest of the report reconstructs the implementation from the source
code, beginning with environment preparation and dataset reconstruction
and then moving through preprocessing, representation learning, anomaly
scoring, calibration, evaluation, robustness analysis, and artifact
packaging. The emphasis is on code-to-concept traceability and on
distinguishing what the program guarantees from what the surrounding
comments merely intend.

\section{2. Project Overview}\label{project-overview}

Conceptually, the implementation consists of the following stages:

\begin{enumerate}
\def\labelenumi{\arabic{enumi}.}
\tightlist
\item
  Prepare a Colab-compatible software environment and reinitialize
  Pillow.
\item
  Create a project directory tree for data, results, predictions, logs,
  models, and checkpoints.
\item
  Fix a small set of seeds and select execution-device parameters.
\item
  Download the compact VisA PCB1 Parquet files from Hugging Face and
  verify their SHA-256 hashes.
\item
  Reconstruct image and mask files from binary Parquet fields and
  generate a standardized \texttt{1cls.csv} index.
\item
  Define deterministic image preprocessing and two families of image
  corruptions.
\item
  Load frozen pretrained ResNet-50 and DINOv2 ViT-S/14 backbones.
\item
  Extract and cache clean embeddings for the full reconstructed dataset.
\item
  For each experimental seed, split normal training data into a
  reference bank and a separate calibration pool.
\item
  Compute k-nearest-neighbor anomaly scores against the reference bank.
\item
  Construct synthetic calibration anomalies from the calibration pool
  and fit logistic and isotonic score calibrators.
\item
  Evaluate raw anomaly scores and calibrated probabilities on untouched
  real VisA test samples.
\item
  Repeat the core experiment across two backbones and three seeds.
\item
  Examine the sensitivity of DINOv2 scoring to reference-bank size and
  the number of neighbors.
\item
  Stress-test the detector under blur, brightness, JPEG compression, and
  resolution degradation at three severity levels.
\item
  Produce reliability and qualitative error-analysis artifacts.
\item
  Package tables, predictions, logs, figures, and the standardized split
  file into a ZIP archive.
\end{enumerate}

The central methodological path can be expressed mathematically as

\[
x \xrightarrow{f_\theta} z \xrightarrow{\ell_2\text{-normalize}} \hat z
\xrightarrow{\text{kNN against }\mathcal{B}} s(x)
\xrightarrow{g} p(x),
\]

where \(x\) is an image, \(f_\theta\) is a frozen backbone, \(\hat z\)
is its normalized embedding, \(\mathcal{B}\) is a reference bank of
normal embeddings, \(s(x)\) is the raw anomaly score, and \(g\) is one
of the fitted calibration functions that maps a scalar score to an
estimated anomaly probability.

The key architectural point is that the code does not learn a feature
extractor and an anomaly decision rule jointly. The backbones
are frozen, and the learned part of the pipeline is limited to
calibration functions applied to the scalar anomaly score. This makes
the experimental question narrower and more interpretable: how useful
are the pretrained representations for nearest-neighbor anomaly
discrimination, and how well can the resulting distance scale be
converted into a probabilistic output under the chosen calibration
protocol?

\section{3. Technical Foundations}\label{technical-foundations}

\subsection{3.1 One-class anomaly detection in representation
space}\label{one-class-anomaly-detection-in-representation-space}

Let \(\mathcal{X}_{\mathrm{normal}}\) denote the distribution of nominal
images. A representation model
\(f_\theta:\mathcal{X}\rightarrow\mathbb{R}^{d}\) maps each image into a
feature vector. Instead of fitting a conventional binary classifier, the
implementation stores a finite reference set

\[
\mathcal{B}=\left\{\hat z_1,\hat z_2,\ldots,\hat z_B\right\},
\]

where each \(\hat z_j\) is a normalized embedding of a normal training
image.

For a query embedding \(\hat z\), small distances to members of
\(\mathcal{B}\) indicate local agreement with known normal appearance.
Large distances indicate that the query is unlike nearby reference
examples and is therefore treated as more anomalous. This is a local
density or neighborhood-consistency view of anomaly detection rather
than a parametric estimate of a complete normal-data distribution.

\subsection{3.2 Frozen transfer
representations}\label{frozen-transfer-representations}

The feature extractors are pretrained models imported from external
libraries. The source code places both models in evaluation mode and
disables gradients for every parameter. Consequently, the experiment is
not a fine-tuning study. It is a frozen-representation study.

The ResNet branch replaces the final fully connected classification
layer with an identity mapping. Therefore, the network output is the
representation immediately before the original ImageNet classifier.

The DINOv2 branch uses \texttt{forward\_features} and constructs

\[
z_{\mathrm{DINO}}
=
\left[
z_{\mathrm{CLS}}
\;\Vert\;
\frac{1}{P}\sum_{p=1}^{P} z_{\mathrm{patch},p}
\right],
\]

where \(z_{\mathrm{CLS}}\) is the class-token representation,
\(z_{\mathrm{patch},p}\) is the representation of patch \(p\), \(P\) is
the number of patch tokens, and \(\Vert\) denotes concatenation. This
should be read as a two-source representation: one global token and one
mean-pooled summary of local patch representations. Although the source
describes this as a multi-scale representation, it is not a conventional
image-pyramid or multi-resolution feature hierarchy; the implementation
uses two token-derived summaries from a single ViT forward pass.

\subsection{3.3 Normalization and cosine
geometry}\label{normalization-and-cosine-geometry}

The code applies

\[
\hat z
=
\frac{z}{\lVert z\rVert_2}
\]

to every extracted embedding. For normalized vectors, the Euclidean
inner product satisfies

\[
\hat z_i^{\top}\hat z_j
=
\cos\!\left(\angle(\hat z_i,\hat z_j)\right).
\]

The implementation therefore defines the pairwise distance as

\[
d(\hat z_i,\hat z_j)
=
1-\hat z_i^{\top}\hat z_j.
\]

This is a cosine-distance formulation on L2-normalized features. It is
especially convenient here because the entire query-to-reference
similarity matrix can be computed as a matrix multiplication.

\subsection{3.4 k-nearest-neighbor anomaly
scoring}\label{k-nearest-neighbor-anomaly-scoring}

For a query embedding \(\hat z\), let
\(d_{(1)}(\hat z),\ldots,d_{(k)}(\hat z)\) be the \(k\) smallest
distances to the reference bank. The anomaly score is

\[
s(\hat z)
=
\frac{1}{k}
\sum_{j=1}^{k}
d_{(j)}(\hat z).
\]

Higher values indicate that even the nearest normal reference examples
are relatively distant.

The use of an average over the nearest \(k\) neighbors makes the score
less sensitive than a one-neighbor rule to one unusually similar or
dissimilar reference image. The code exposes \(k\) as a hyperparameter
and uses \(k=5\) for the main experiment.

\subsection{3.5 Calibration as a separate probabilistic
layer}\label{calibration-as-a-separate-probabilistic-layer}

The raw kNN score is a distance-derived quantity, not a probability. The
code therefore fits two one-dimensional calibration models.

The logistic model estimates

\[
p_{\mathrm{log}}(y=1\mid s)
=
\sigma(as+b)
=
\frac{1}{1+\exp[-(as+b)]},
\]

where \(a\) and \(b\) are fitted parameters and \(\sigma\) is the
logistic sigmoid.

The isotonic model instead learns a non-decreasing function

\[
p_{\mathrm{iso}}(y=1\mid s)=g(s),
\]

subject to the monotonicity constraint

\[
s_1\le s_2
\quad\Rightarrow\quad
g(s_1)\le g(s_2).
\]

Monotonicity is natural here because the anomaly score is intended to
increase with abnormality.

\subsection{3.6 Calibration metrics}\label{calibration-metrics}

The implementation evaluates calibration with adaptive expected
calibration error, the Brier score, and a risk-coverage calculation.

For adaptive ECE, predictions are partitioned using empirical quantiles
of the predicted probabilities. For bin \(b\), let \(n_b\) be its sample
count, \(c_b\) the mean predicted probability, and \(f_b\) the empirical
anomaly frequency. The reported quantity is

\[
\operatorname{ECE}
=
\sum_{b=1}^{M}
\frac{n_b}{N}
\left|c_b-f_b\right|.
\]

The use of quantile-based bins is intended to reduce empty or extremely
sparse bins when the probability distribution is highly concentrated.
The implementation falls back to equal-width bins if the quantile
construction collapses to too few unique boundaries.

The Brier score is

\[
\operatorname{BS}
=
\frac{1}{N}
\sum_{i=1}^{N}
\left(p_i-y_i\right)^2.
\]

Lower values indicate that predicted probabilities are, on average,
closer to the observed binary outcomes.

For risk-coverage analysis, the predicted class is

\[
\hat y_i=\mathbf{1}[p_i\ge 0.5],
\]

and confidence is

\[
c_i=\max(p_i,1-p_i).
\]

Samples are sorted from highest to lowest confidence. At each coverage
level, the cumulative classification error defines empirical risk. The
code approximates the area under this risk-coverage curve using
numerical trapezoidal integration.

\subsection{3.7 Discrimination metrics}\label{discrimination-metrics}

The code uses AUROC and average precision (AUPRC) on the raw anomaly
scores. Because larger scores correspond to stronger anomaly evidence, a
higher score is treated as a positive anomaly ranking.

AUROC measures ranking quality across possible score thresholds. Average
precision summarizes the precision-recall relationship and is sensitive
to class prevalence. Neither metric directly establishes probabilistic
calibration.

\section{4. Input Data and
Reconstruction}\label{input-data-and-reconstruction}

\subsection{4.1 Dataset source and integrity
checks}\label{dataset-source-and-integrity-checks}

The source obtains compact Parquet files from the Hugging Face dataset
repository identified in the code as \texttt{BrachioLab/visa}. Separate
files are used for training and testing. Before downloading a file, the
implementation checks whether a local copy exists and whether its
SHA-256 digest matches a hard-coded expected value. If the local digest
does not match, the file is deleted and redownloaded.

The digest function reads each file in chunks and updates a SHA-256 hash
object incrementally:

\[
H
=
\operatorname{SHA256}
\left(
\text{bytes of the complete file}
\right).
\]

This is useful for reproducibility because it checks file identity
rather than merely file size or presence.

The implementation copies the verified cached download into the
experiment's own data directory. The source of the file is therefore
explicit, and the local copy is independently checked after copying.

\subsection{4.2 Binary image decoding}\label{binary-image-decoding}

The Parquet records contain image and mask fields whose exact Python
representation may be bytes-like objects or dictionaries containing byte
payloads. \texttt{\_bytes\_from\_value} normalizes these cases into a
\texttt{bytes} value. Unsupported representations raise a
\texttt{TypeError}, which is preferable to silently treating unexpected
data structures as valid images.

RGB images are decoded with Pillow, explicitly converted to RGB, and
saved as JPEG with quality 95. Masks are decoded as grayscale images and
stored as PNG files. The masks are stored only for anomalous samples.

The reconstruction step queries the Parquet files through DuckDB:

\[
\texttt{SELECT image, mask, label FROM read\_parquet(...)}.
\]

Each returned record is assigned a deterministic file name containing
the split name and row index. The binary class rule is explicitly

\[
\text{is\_anomaly} = [\text{label\_id}=1].
\]

Thus, the implementation assumes the source labels use \texttt{1} for
anomaly and treats the alternative label value as normal.

\subsection{4.3 Standardized split
index}\label{standardized-split-index}

The reconstruction process produces a row-level table with object name,
split, string label, image path, mask path, and integer label. The table
is written to \texttt{1cls.csv}, providing a stable index for subsequent
processing.

The code subsequently filters this table rather than re-reading the raw
Parquet files. In the main experiment, only images with
\texttt{split\ ==\ \textquotesingle{}train\textquotesingle{}} and
\texttt{label\ ==\ \textquotesingle{}normal\textquotesingle{}} are
eligible for the reference and calibration pools. The test frame
includes all test images, including real anomalies.

This separation is central to the leakage-free design. The real test
labels are used only for evaluation, after the reference bank and
calibrators have been fixed.

\section{5. Preprocessing and Synthetic
Corruptions}\label{preprocessing-and-synthetic-corruptions}

\subsection{5.1 Image preprocessing}\label{image-preprocessing}

The common preprocessing pipeline performs three operations:

\[
x
\rightarrow
\operatorname{Resize}(x,224,224)
\rightarrow
\operatorname{ToTensor}(x)
\rightarrow
\operatorname{Normalize}(x;\mu,\sigma),
\]

with ImageNet channel statistics

\[
\mu=(0.485,0.456,0.406),
\qquad
\sigma=(0.229,0.224,0.225).
\]

The fixed square resize means that original aspect ratio is not
preserved. The code does not apply a center crop, random crop, or
explicit geometric normalization. Consequently, all inputs are forced to
the same spatial resolution before entering either pretrained backbone.

\subsection{5.2 Synthetic defects for
calibration}\label{synthetic-defects-for-calibration}

The calibration stage is unusual in an important way: it needs positive
anomaly examples even though the true anomalies in the held-out test set
cannot be used for fitting. The implementation addresses this by
perturbing normal calibration images with controlled synthetic defects.

Four calibration defect families are implemented.

\subsubsection{Subtle additive noise}\label{subtle-additive-noise}

Pixel values are converted to the interval \([0,1]\). A random Gaussian
perturbation with standard deviation sampled uniformly from \(0.015\) to
\(0.035\) is added and clipped to \([0,1]\):

\[
x'=\operatorname{clip}(x+\epsilon,0,1),
\qquad
\epsilon\sim\mathcal{N}(0,\sigma^2 I).
\]

The code samples a new \(\sigma\) for each image.

\subsubsection{Cut-and-paste patch}\label{cut-and-paste-patch}

A patch with width and height independently sampled from the integer
range 12 to 27 pixels is taken from one image location and copied to
another. With probability one half, the patch is rotated by
\(180^\circ\) before insertion. This creates a localized appearance
inconsistency while preserving most of the image.

\subsubsection{Subtle blur}\label{subtle-blur}

A Gaussian blur radius is sampled uniformly from \(0.4\) to \(0.9\).
This introduces a localized-defect calibration example only in the broad
sense that the transformed image is presented as a positive synthetic
anomaly; spatial localization is not explicitly controlled by this
corruption.

\subsubsection{Micro-scratch}\label{micro-scratch}

A short line is generated by selecting a starting point, a random angle,
and a length between 15 and 39 pixels. The affected pixels are
brightened by a random amount between 60 and 119 intensity units before
clipping.

The docstring for the synthetic-defect function states that the
resulting calibration scores are guaranteed to occupy a particular
numerical range and match the scale of real defects. The implementation
itself does not enforce such a score range. The scores emerge from the
backbone and reference bank. Therefore, this statement should be
interpreted as an intended design claim, not as a mathematical
guarantee.

\subsection{5.3 Stress-test corruptions}\label{stress-test-corruptions}

A separate corruption function is used for out-of-distribution stress
testing. It supports:

\begin{itemize}
\tightlist
\item
  Gaussian blur with radius \(0.7s\), where \(s\in\{1,2,3\}\) is
  severity.
\item
  Brightness scaling using factors \(1.15\), \(0.75\), and \(0.50\).
\item
  JPEG recompression with qualities 75, 45, and 20.
\item
  Resolution degradation by downscaling to \(65\%\), \(40\%\), or
  \(25\%\) of the original dimensions and resizing back to the original
  size.
\end{itemize}

These transformations are deterministic with respect to the image-level
seed and severity parameters used by the dataset wrapper.

The conceptual distinction between the two corruption families matters.
Calibration corruptions are training-pool transformations used to fit
the probability mapping, while stress corruptions are applied to the
untouched test set after the calibrators have been fit. The latter
therefore test behavior under distribution shift rather than augmenting
the fitting process.

\section{6. Dataset Abstraction and Execution
Model}\label{dataset-abstraction-and-execution-model}

\texttt{PathImageDataset} implements the PyTorch \texttt{Dataset}
interface. It stores a copy of the incoming dataframe, the common
transform, optional corruption settings, and a base seed.

For an index \(i\), \texttt{\_\_getitem\_\_} performs:

\[
\begin{aligned}
\text{row}_i
&\rightarrow \text{read image}
\rightarrow \text{RGB conversion}
\rightarrow \text{optional corruption}\\
&\rightarrow \text{tensor transform}
\rightarrow (\text{image tensor},\text{label},\text{path}).
\end{aligned}
\]

The dataset passes a seed of the form

\[
s_i=s_{\mathrm{base}}+i
\]

to corruption routines. This makes the image-level transformation
deterministic for a fixed dataset ordering and seed.

The \texttt{DataLoader} uses \texttt{shuffle=False}, which is important
because the returned paths are later used to map embeddings back to
dataframe rows. The loader also uses pinned memory when the GPU is
active and chooses the worker count based on the execution device.

One subtle design decision is that feature extraction is path-indexed
after the fact. The code extracts embeddings in loader order, saves the
paths alongside those embeddings, and later constructs a dictionary from
path to feature-row index. This avoids assuming that dataframe row order
and feature-array order are permanently interchangeable.

\section{7. Backbone Loading and Feature
Extraction}\label{backbone-loading-and-feature-extraction}

\subsection{7.1 ResNet-50 branch}\label{resnet-50-branch}

The ResNet-50 model is loaded with the library's default pretrained
weights. Its final classification layer is replaced with
\texttt{nn.Identity()}:

\[
f_{\theta}^{\mathrm{ResNet}}(x)
=
h_{\theta}(x),
\]

where \(h_\theta(x)\) is the feature vector immediately preceding the
original classifier.

The model is placed on the selected device and switched to evaluation
mode. Every parameter is marked \texttt{requires\_grad=False}, so no
optimization graph is constructed for the model.

This is a conventional way to repurpose a classification network as a
fixed representation function.

\subsection{7.2 DINOv2 branch}\label{dinov2-branch}

DINOv2 is loaded through \texttt{torch.hub} using the
\texttt{dinov2\_vits14} entry point. The same frozen-model treatment is
applied.

The source explicitly accesses the dictionary returned by
\texttt{forward\_features} and reads the normalized class token and
normalized patch tokens. Patch tokens are averaged along the patch
dimension:

\[
\bar z_{\mathrm{patch}}
=
\frac{1}{P}\sum_{p=1}^{P}z_{\mathrm{patch},p},
\]

then concatenated with the class token.

The code therefore creates a representation that combines a global
semantic token with the average of local patch-level representations. The patch mean is computationally simple, but it
deliberately discards the spatial arrangement among patches after
aggregation.

\subsection{7.3 Mixed-precision
inference}\label{mixed-precision-inference}

When CUDA is active, the source wraps the forward pass in
\texttt{torch.amp.autocast(\textquotesingle{}cuda\textquotesingle{})}.
CPU execution uses a null context.

After inference, embeddings are converted to float32 and normalized
before they are copied to CPU NumPy arrays. This arrangement reduces
intermediate GPU memory pressure while ensuring a consistent float32
representation for the downstream NumPy calculations.

\subsection{7.4 Feature caching}\label{feature-caching}

Clean features are cached as a compressed \texttt{.npz} file together
with a CSV containing the corresponding paths and labels.

Caching is beneficial because the same frozen representations are reused
across multiple seeds, calibration trials, hyperparameter checks, and
robustness experiments. The cache is particularly valuable when the
backbones are the computational bottleneck.

However, the cache key is effectively just the backbone name. The cache
does not record the source-image hash, preprocessing configuration,
model weight identity, library version, or feature dimensionality. A
stale cache could therefore be reused after a change in model or
preprocessing configuration if the expected files happen to remain
present. This is an implementation-level reproducibility limitation.

\section{8. Leakage-Free Reference and Calibration
Protocol}\label{leakage-free-reference-and-calibration-protocol}

\subsection{8.1 Category filtering}\label{category-filtering}

\texttt{get\_category\_frames} constructs two frames for the selected
object category. The training frame contains only normal examples; the
test frame contains all test examples.

In this configuration, the category list contains only
\texttt{pcb1}. Thus, although the helper functions accept a category
argument, the actual experiment is single-category rather than a
multi-category benchmark.

\subsection{8.2 Reference and calibration
split}\label{reference-and-calibration-split}

The normal training frame is randomly permuted using a NumPy generator
seeded with the trial seed. A maximum of 250 samples is assigned to the
reference bank and a further maximum of 150 samples is assigned to the
calibration pool.

The split is without replacement because the calibration indices are
taken from the unused portion after the bank indices have been selected.

Formally, for training normals \(\mathcal{T}\),

\[
\mathcal{T}
=
\mathcal{B}_{\mathrm{data}}
\cup
\mathcal{C}_{\mathrm{data}}
\cup
\mathcal{R},
\]

where \(\mathcal{B}_{\mathrm{data}}\) is the reference subset,
\(\mathcal{C}_{\mathrm{data}}\) is the calibration subset, and
\(\mathcal{R}\) is the remainder not used by the main method.

The separation is scientifically important. If the same examples defined
both the normal reference bank and the calibration mapping, the
calibration score distribution would be tied to the exact points used to
define normality. The code avoids that coupling by construction.

\subsection{8.3 Converting dataframe rows to feature
rows}\label{converting-dataframe-rows-to-feature-rows}

The feature cache contains an array and a parallel list of paths.
\texttt{frame\_to\_feature\_rows} constructs a path-to-index dictionary
and retrieves the correct feature rows for a selected dataframe.

This path-based indexing is an explicit data-integrity safeguard. It
reduces the chance that later dataframe filtering or concatenation
silently changes the association between metadata and embeddings.

\section{9. Anomaly Scoring and
Calibration}\label{anomaly-scoring-and-calibration}

\subsection{9.1 Chunked similarity
computation}\label{chunked-similarity-computation}

\texttt{score\_knn} processes query embeddings in chunks. For a query
block \(Q\) and reference matrix \(R\), it computes

\[
S=QR^\top,
\]

where each matrix element is the inner product between one normalized
query and one normalized reference vector. Because the vectors are
normalized,

\[
D=\mathbf{1}-S
\]

is the corresponding cosine-distance matrix.

For every query row, the implementation uses \texttt{np.partition} to
obtain the \(k\) smallest distances without fully sorting all reference
distances. It then averages those nearest values.

This is efficient for the intended use case. A full sort would
perform more work than necessary when only the first \(k\) elements are
required.

\subsection{9.2 Calibration data
construction}\label{calibration-data-construction}

For a trial, the calibration score set consists of:

\begin{enumerate}
\def\labelenumi{\arabic{enumi}.}
\tightlist
\item
  scores from clean normal calibration images, labeled \(0\);
\item
  scores from cut-and-paste calibration variants, labeled \(1\);
\item
  scores from subtle-noise variants, labeled \(1\);
\item
  scores from subtle-blur variants, labeled \(1\);
\item
  scores from micro-scratch variants, labeled \(1\).
\end{enumerate}

With the configured calibration-pool size of 150, the intended
construction contains up to 150 clean examples and up to \(4\times150\)
synthetic positive examples, provided the source pool is large enough.

The key methodological property is not the exact sample count but the
information boundary: real test images never contribute to the
calibration fit.

\subsection{9.3 Logistic calibration}\label{logistic-calibration}

The logistic regression uses a single predictor: the scalar anomaly
score. Its fitted decision function can be written as

\[
\eta(s)=as+b,
\]

with probability

\[
p_{\mathrm{log}}(s)=\sigma(\eta(s)).
\]

The regularization parameter is explicitly set to \(C=1.0\), the solver
is L-BFGS, and the maximum iteration count is 1000.

The model is therefore intentionally low-dimensional. It does not learn
a representation; it learns only a monotone sigmoid-like remapping of
score magnitude.

\subsection{9.4 Isotonic calibration}\label{isotonic-calibration}

The isotonic regressor fits a monotone function directly to the
calibration pairs \((s_i,y_i)\). Its output is constrained to the
interval \([0,1]\) and is clipped for out-of-bounds scores.

Isotonic calibration is more flexible than logistic calibration because
it does not impose a particular sigmoid shape. At the same time, its
flexibility can make it sensitive to calibration-set composition,
especially when the calibration sample is not large.

\section{10. Evaluation Procedure}\label{evaluation-procedure}

\subsection{10.1 Main trial}\label{main-trial}

The core experiment is repeated over

\[
3\ \text{seeds}\times 2\ \text{backbones}\times 1\ \text{category},
\]

giving six main trials.

For each trial, the reference bank and calibration pool are resampled
from normal training data according to the trial seed. The calibrators
are fitted only on that trial's calibration data.

The real test images are then scored against the reference bank. AUROC
and AUPRC are computed from the raw anomaly scores, while ECE, Brier
score, and AURC are computed separately for logistic and isotonic
probabilities.

The raw score is therefore the object of discrimination evaluation, and
the calibrated probability is the object of probabilistic evaluation.

\subsection{10.2 What is and is not
aggregated}\label{what-is-and-is-not-aggregated}

The source stores one result row per backbone-category-seed combination
and computes means and standard deviations grouped by backbone for
selected metrics. These summaries are aggregation operations, not
uncertainty estimates in the inferential-statistical sense. The code
does not compute confidence intervals, hypothesis tests, bootstrap
intervals, or repeated nested resampling.

The resulting standard deviation across three seeds should therefore be
interpreted as a small-sample estimate of variability under the selected
random-split protocol, not as a formal population-level uncertainty
bound.

\subsection{10.3 Hyperparameter
sensitivity}\label{hyperparameter-sensitivity}

The reference-bank size \(B\) is varied over

\[
B\in\{50,150,250\},
\]

and the neighbor count \(k\) is varied over

\[
k\in\{1,5,10\}.
\]

This creates nine configurations evaluated under three seeds, using
DINOv2 features.

The sensitivity study is informative about how the kNN detector responds
to memory budget and neighborhood size. However, the same real test set
is used to compare all configurations. Therefore, this analysis should
be viewed as an evaluation or sensitivity study unless an external
decision rule has already specified the final hyperparameters. It should
not be described as a fully nested hyperparameter-selection procedure.

\section{11. Robustness Stress-Testing Under Distribution
Shift}\label{robustness-stress-testing-under-distribution-shift}

The robustness stage fixes the reference/calibration split to seed 13
and fits calibrators in the same leakage-free manner as the main
experiment. It then evaluates the untouched test set under both clean
and corrupted conditions.

For each backbone, the stress suite covers four corruption types and
three severities, in addition to a clean condition.

The raw anomaly detector is kept unchanged. Only the test images are
transformed. Consequently, the experiment probes whether the
representation and kNN score maintain useful discrimination when image
appearance changes while the underlying normal/anomalous label remains
fixed.

The calibrated probability is also passed through the same logistic or
isotonic mapping learned from clean plus subtle synthetic calibration
data. This is important: the stress experiment does not re-fit the
calibrators under the corruption. Calibration degradation is therefore a
measure of how the previously learned probability mapping transfers to
shifted score distributions.

This protocol separates two failure modes:

\[
\text{distribution shift}
\rightarrow
\begin{cases}
\text{ranking degradation},\\
\text{probability-calibration degradation}.
\end{cases}
\]

A model may preserve AUROC while becoming poorly calibrated, or it may
lose both ranking and calibration quality. The source explicitly records
both families of measures.

\section{12. Reliability Analysis and Qualitative Error
Diagnosis}\label{reliability-analysis-and-qualitative-error-diagnosis}

\subsection{12.1 Reliability diagram}\label{reliability-diagram}

The reliability-diagram function evaluates the DINOv2 seed-13
predictions using adaptive calibration bins. It compares mean predicted
probability with empirical anomaly frequency.

Although the source generates a graphical reliability diagram and
histogram, this report deliberately does not reproduce those figures.
The underlying computational logic is still important because it
provides a visual companion to ECE: ECE is a single scalar, whereas a
reliability relationship reveals where along the probability axis the
model is over- or under-confident.

\subsection{12.2 Ablation table}\label{ablation-table}

The ablation table records three conceptual variants per backbone:

\[
\text{raw score},
\qquad
\text{score + logistic calibration},
\qquad
\text{score + isotonic calibration}.
\]

The AUROC and AUPRC values are copied unchanged across the three
variants because calibration is applied after the raw ranking has been
produced. ECE and Brier score are only reported for the calibrated
variants.

This separates ranking effects from calibration effects. A monotone
calibration transformation can change the numerical interpretation of a score without changing its ordering,
so ranking metrics should remain unchanged unless numerical ties or
implementation details intervene.

\subsection{12.3 False-positive and false-negative
selection}\label{false-positive-and-false-negative-selection}

The qualitative error-analysis code focuses on the DINOv2 seed-13 trial.
It selects the four highest-scoring normal test examples as false
positives and the four lowest-scoring real anomalies as false negatives.

The selection rule is direct:

\[
\text{FP}
=
\operatorname{Top4}
\left\{
s_i:y_i=0
\right\},
\]

\[
\text{FN}
=
\operatorname{Bottom4}
\left\{
s_i:y_i=1
\right\}.
\]

The source then saves image grids with raw scores and isotonic
probabilities. Those images can be useful for expert diagnosis of
representation failure modes, but no automated semantic interpretation
of the failure cases is performed in the code. Therefore, the report
does not assign visual causes to individual errors.

\section{13. Software Architecture and Code
Organization}\label{software-architecture-and-code-organization}

The program is organized as a linear research script with helper
functions rather than as a multi-module Python package. This is
appropriate for interactive Colab execution but has consequences for
maintainability.

The code can be understood as nine functional layers:

\begin{enumerate}
\def\labelenumi{\arabic{enumi}.}
\tightlist
\item
  \textbf{Environment reset:} module purging, Pillow replacement,
  runtime restart.
\item
  \textbf{Configuration:} directory creation, seeds, device, backbone
  names, batch size, image size.
\item
  \textbf{Dataset reconstruction:} download, hash verification,
  image/mask materialization, CSV generation.
\item
  \textbf{Synthetic data generation:} calibration and stress
  corruptions.
\item
  \textbf{Dataset abstraction:} \texttt{PathImageDataset}.
\item
  \textbf{Representation extraction:} backbone loading, inference,
  feature normalization, caching.
\item
  \textbf{Scoring and calibration:} reference sampling, kNN score,
  calibration models, metrics.
\item
  \textbf{Experiment orchestration:} main trials, sensitivity analysis,
  robustness analysis, qualitative diagnostics.
\item
  \textbf{Packaging:} result-directory assembly and ZIP creation.
\end{enumerate}

The principal state is kept in module-level variables such as
\texttt{split\_df}, \texttt{clean\_features}, \texttt{MODELS}, and the
result dataframes. This makes the script easy to execute top-to-bottom
in a notebook-like environment, but it also means individual functions
depend on global state rather than explicit dependency injection.

\section{14. Detailed Control Flow}\label{detailed-control-flow}

The intended top-to-bottom execution is important because the file is
designed for Colab.

The environment cell first detects whether Colab is present. In Colab,
the process exits intentionally after reinstalling Pillow. The notebook
must then be rerun from the beginning so that the new C-extension is
loaded into a fresh process.

Once restarted, the script defines the project directories and imports
all computational libraries. It then reconstructs the dataset and writes
the split table.

The backbones are loaded only after the dataset abstraction is defined.
Clean embeddings are extracted once per backbone and cached. The main
experimental loop then reuses those embeddings rather than repeatedly
running the frozen networks on clean images.

Synthetic calibration features are extracted inside each trial because
the transformed images differ by corruption type and seed. Stress-test
features are also extracted on demand because they depend on corruption
and severity.

Finally, the script writes tables and predictions, generates diagnostic
graphics, and packages the artifacts. The \texttt{final\_package}
directory is rebuilt on every run by deleting any previous directory and
recreating it.

This execution order matches the computational dependencies: model
loading precedes feature extraction, and feature caching precedes
repeated scoring analyses.

\section{15. Design Decisions and
Rationale}\label{design-decisions-and-rationale}

The rationale below follows from the implementation choices, although the
author's intent is not separately documented.

\subsection{15.1 Frozen backbones rather than task-specific
training}\label{frozen-backbones-rather-than-task-specific-training}

Freezing ResNet-50 and DINOv2 isolates the effect of representation
quality from the optimization behavior of a trainable anomaly model. It
also reduces the risk of overfitting a small industrial dataset to a
particular defect taxonomy.

The trade-off is that no feature adaptation is possible. If the
pretrained representation does not encode the appearance differences
that matter for PCB defects, the downstream kNN score cannot recover
that information.

\subsection{15.2 Reference-bank scoring rather than a parametric density
model}\label{reference-bank-scoring-rather-than-a-parametric-density-model}

The kNN score requires comparatively few assumptions about the shape of
normal embeddings. It does not assume Gaussian features, a specific
covariance structure, or an explicit probability density.

Its main computational cost grows with the reference-bank size, which
motivates the explicit budget experiment.

\subsection{15.3 Separate calibration
pool}\label{separate-calibration-pool}

The separate calibration pool is one of the strongest experimental
design features in the source. It avoids fitting the
score-to-probability function on the same normal images used to define
the reference neighborhood.

\subsection{15.4 Synthetic anomalies only for
calibration}\label{synthetic-anomalies-only-for-calibration}

This permits calibration without consuming real test anomalies. The
method is pragmatic, but synthetic calibration examples need not have
the same score distribution as real industrial defects. The calibration
result is therefore dependent on the realism and coverage of the chosen
corruption family.

\subsection{15.5 Two calibrators}\label{two-calibrators}

Logistic regression supplies a simple parametric baseline. Isotonic
regression supplies a flexible monotone alternative. Together they allow
the code to distinguish the question ``does a simple sigmoid remapping
suffice?'' from the question ``does a non-parametric monotone mapping
better fit the observed calibration relationship?''

\subsection{15.6 Stress-testing after
calibration}\label{stress-testing-after-calibration}

Applying distribution-shift corruptions after the calibrators are fixed
is methodologically appropriate for evaluating transfer of calibration.
Recalibrating on shifted data would answer a different question.

\section{16. Mathematical and Computational
Complexity}\label{mathematical-and-computational-complexity}

\subsection{16.1 Feature extraction}\label{feature-extraction}

Let \(N\) be the number of images, and let \(C_f\) denote the
computational cost of one backbone forward pass. Feature extraction
scales approximately as

\[
O(N\,C_f),
\]

with the exact constant dominated by the selected backbone and image
size.

The use of batching reduces Python-level overhead and can improve GPU
utilization. Mixed precision on CUDA can further reduce memory traffic
and improve throughput, depending on hardware and PyTorch execution.

\subsection{16.2 kNN scoring}\label{knn-scoring}

Let \(Q\) be the number of query embeddings, \(B\) the reference-bank
size, and \(d\) the embedding dimension. Computing the similarity matrix
costs approximately

\[
O(QBd).
\]

Finding the smallest \(k\) distances with \texttt{np.partition} avoids
an \(O(B\log B)\) full sort for every query row. The exact complexity of
the selection primitive depends on the NumPy implementation, but the key
design point is that only the nearest \(k\) values are required.

Peak temporary memory is controlled by the chunk size \(C\), giving a
similarity matrix of roughly \(C\times B\) rather than \(Q\times B\).
The corresponding temporary memory therefore scales as

\[
O(CB).
\]

\subsection{16.3 Cache and storage
trade-offs}\label{cache-and-storage-trade-offs}

The feature cache trades disk space for repeated computation. This is
appropriate because the clean embeddings are reused many times, while
the backbones are the most expensive part of the pipeline.

The reconstruction step also expands compact Parquet records into
individual JPEG and PNG files. That simplifies downstream image loading
and standardizes file-based access, but it increases the number of
filesystem objects and introduces a JPEG recompression step.

\section{17. Reproducibility and
Reliability}\label{reproducibility-and-reliability}

\subsection{17.1 Positive reproducibility
mechanisms}\label{positive-reproducibility-mechanisms}

The implementation includes several useful reproducibility practices:

\begin{itemize}
\tightlist
\item
  explicit seed values are defined centrally;
\item
  NumPy random generators are used for dataset splits and corruptions;
\item
  the exact image-transform parameters are in source;
\item
  dataset files are checked with SHA-256 digests;
\item
  clean features are cached with explicit path metadata;
\item
  experiment tables and predictions are written to deterministic
  locations;
\item
  a standardized split CSV is created from the reconstructed dataset.
\end{itemize}

These practices improve repeatability and auditability.

\subsection{17.2 Why reproducibility is not fully
deterministic}\label{why-reproducibility-is-not-fully-deterministic}

The source calls

\begin{Shaded}
\begin{Highlighting}[]
\NormalTok{torch.backends.cudnn.benchmark }\OperatorTok{=} \VariableTok{True}
\end{Highlighting}
\end{Shaded}

when initializing seeds. cuDNN benchmarking can select algorithms based
on runtime characteristics and is not the same as requesting
deterministic kernels. The source does not set
\texttt{torch.backends.cudnn.deterministic\ =\ True} and does not
globally disable nondeterministic operations.

Thus, the correct characterization is \textbf{seeded and partially
controlled}, not strictly bitwise deterministic.

The script also uses multiprocessing in the \texttt{DataLoader} when
CUDA is available. Although the corruption itself is deterministically
parameterized by an image index and base seed, exact end-to-end
reproducibility can still depend on worker behavior and library
implementations.

\subsection{17.3 External dependency
reproducibility}\label{external-dependency-reproducibility}

Pillow is explicitly pinned to \texttt{12.3.0}. In contrast, the source
does not pin the versions of PyTorch, torchvision, DuckDB, scikit-learn,
NumPy, pandas, or the Hugging Face tooling. Their versions are printed
at runtime, but printing a version does not by itself recreate an
environment.

The DINOv2 weights are loaded through \texttt{torch.hub}, and the code
does not record a repository commit or model-weight checksum. ResNet
uses the library's \texttt{DEFAULT} weight specification rather than an
explicit immutable artifact identifier.

These choices are common in exploratory Colab workflows, but they limit
long-term reproducibility.

\subsection{17.4 Environment reset as a reliability
measure}\label{environment-reset-as-a-reliability-measure}

The Pillow reset logic is unusually aggressive: it purges loaded PIL
modules, uninstalls Pillow, manually removes matching site-package
directories, reinstalls a chosen version, and terminates the Colab
process.

The intent is to ensure that compiled Pillow extensions come from a
clean installation. The approach can resolve stale binary-extension
problems, but it is also destructive and tightly coupled to a mutable
Colab environment. It is not a portable environment-management strategy
for ordinary local execution.

\section{18. Limitations and Technical
Considerations}\label{limitations-and-technical-considerations}

\subsection{18.1 Calibration prevalence is
synthetic}\label{calibration-prevalence-is-synthetic}

The calibration labels are constructed from one clean set and four
synthetic anomaly sets. With a calibration pool of 150 images, this can
create a nominal calibration class ratio of approximately 1 clean
example to 4 synthetic anomaly examples.

That class prevalence is a property of the calibration protocol, not
necessarily the real prevalence of defects in deployment. Because the
logistic intercept depends on class balance, the calibrated
probabilities can inherit a prior-distribution mismatch. Isotonic
calibration can also be affected by the calibration score distribution
and label frequencies.

Therefore, the resulting probabilities should be interpreted as
probabilities under the implemented calibration sample construction
rather than as a universally valid estimate of deployment prevalence.

\subsection{18.2 Synthetic anomalies may not match real
defects}\label{synthetic-anomalies-may-not-match-real-defects}

Cut-and-paste, noise, blur, and scratch transformations are useful
stressors, but they are only proxies for the unknown distribution of
real industrial defects. Calibration quality on these synthetic
positives therefore depends on whether the synthetic examples occupy a
score regime representative of real anomalies.

The code never demonstrates that equivalence. The report consequently
treats synthetic calibration as a methodological device rather than as
evidence that the synthetic and real defect distributions coincide.

\subsection{\texorpdfstring{18.3 \texttt{ece15} naming
mismatch}{18.3 ece15 naming mismatch}}\label{ece15-naming-mismatch}

The result columns are named \texttt{logistic\_ece15} and
\texttt{isotonic\_ece15}. However, the metric function is called with
its default \texttt{n\_bins=10}. The implementation therefore computes a
ten-bin adaptive ECE while naming the output as though it were an ECE-15
metric.

This is a concrete reporting inconsistency. It does not necessarily make
the numerical ECE implementation wrong, but it can create confusion when
comparing the saved table against documentation or another experiment.

\subsection{18.4 Main experiments and secondary analyses use different
seed
coverage}\label{main-experiments-and-secondary-analyses-use-different-seed-coverage}

The primary experiment runs three seeds. The robustness stress test,
reliability diagram, and qualitative error diagnosis use seed 13 only.

These are therefore single-seed diagnostics rather than three-seed
robustness summaries. Their conclusions should not be generalized to the
full seed distribution without additional runs.

\subsection{18.5 Hyperparameter study is not
nested}\label{hyperparameter-study-is-not-nested}

The \(B\)-versus-\(k\) sensitivity study evaluates multiple
configurations against the same real test set. If a configuration were
selected on the basis of those test metrics and then presented as the
final model, the test set would have served as a tuning resource.

The clean interpretation is a post hoc sensitivity analysis that
illustrates how the detector behaves under alternative
parameterizations.

\subsection{18.6 Feature cache invalidation is
implicit}\label{feature-cache-invalidation-is-implicit}

The feature-cache files are named by backbone only. Changes in
preprocessing, source data, model weight version, or embedding
construction could therefore leave a cache file that is technically
readable but scientifically stale.

A stronger implementation would store a manifest containing data hashes,
preprocessing parameters, model identifiers, feature dimensionality, and
perhaps a source-code or configuration hash.

\subsection{18.7 Image reconstruction introduces JPEG
transformation}\label{image-reconstruction-introduces-jpeg-transformation}

The raw RGB images are saved as JPEG with quality 95. If the original
Parquet data are not already encoded in exactly the same way, this step
can change pixel values. The experiment therefore operates on
reconstructed JPEG files rather than directly on the original byte
representation.

For a representation-based detector this may be acceptable, but it is
still a transformation that should be included in a strict
reproducibility description.

\subsection{18.8 Masks are reconstructed but not used in the
detector}\label{masks-are-reconstructed-but-not-used-in-the-detector}

Anomaly masks are extracted and saved for anomalous images, but the
actual scoring, calibration, robustness evaluation, and metrics are
image-level. No pixel-level localization metric, mask overlap, or
segmentation loss is computed.

The presence of masks in the reconstructed dataset should therefore not
be interpreted as evidence that the method performs defect localization.

\subsection{18.9 Multi-scale terminology should be interpreted
carefully}\label{multi-scale-terminology-should-be-interpreted-carefully}

The DINOv2 representation combines a CLS token with a mean of patch
tokens. This is richer than using a single token alone, but it is not a
conventional multi-scale feature pyramid. The code does not extract
feature maps at multiple spatial resolutions.

Calling the method ``multi-scale'' is therefore reasonable only in the
broad sense of combining global and aggregated local token information.

\subsection{18.10 Packaging message is stronger than the implemented
verification}\label{packaging-message-is-stronger-than-the-implemented-verification}

The packaging code writes a \texttt{README.txt} stating that all
verification checks passed. The shown code verifies that the files can
be copied and zipped and reports successful archive creation. It does
not implement a comprehensive semantic verification of every table,
prediction file, or figure.

The message should therefore be understood as a packaging-success
message rather than as evidence of a formal end-to-end artifact audit.

\section{19. Scientific Interpretation
Boundaries}\label{scientific-interpretation-boundaries}

A rigorous reading of the implementation requires separating several
claims that are easy to conflate.

First, the code \textbf{does implement} a leakage-aware protocol in
which the real test anomalies are withheld from the calibration fit.
That is directly supported by the sequence of data selection,
calibration fitting, and final test evaluation.

Second, the code \textbf{does implement} score calibration. It does not
prove that the resulting probabilities are calibrated under any
deployment distribution beyond the empirical conditions represented by
the calibration data.

Third, the code \textbf{does implement} distribution-shift stress
testing. It does not establish invariance to arbitrary industrial
shifts, because the shift family is limited to four hand-designed image
corruptions at three severities.

Fourth, the source \textbf{does compute} several recognized evaluation
metrics. The metrics alone do not establish statistical significance,
generalization to other categories, or superiority over alternative
methods because the shown implementation does not provide a complete
comparative benchmark design.

Finally, the source \textbf{contains diagnostic plots and result
tables}, but this report does not reproduce or interpret their empirical
values. No numerical performance claim is made here beyond what can be
established from code structure and configuration.

\section{20. Overall Technical
Assessment}\label{overall-technical-assessment}

The implementation is technically coherent as a research prototype for
image-level anomaly detection using frozen representations. Its
strongest methodological feature is the explicit separation of reference
construction, calibration, and final evaluation. The code also
demonstrates a useful awareness of practical reproducibility issues
through file hashing, controlled split sampling, deterministic
corruption parameters, feature caching, and explicit artifact packaging.

The choice of cosine-normalized features plus kNN scoring is
computationally straightforward and conceptually transparent. It makes
the anomaly score easy to inspect and separates representation
extraction from anomaly scoring. This separation is especially valuable
for a research-oriented study because changes to the backbone, reference
size, or calibration layer can be analyzed independently.

The calibration design is also methodologically disciplined. Logistic
and isotonic regression are both applied to the same scalar score,
making their comparison relatively interpretable. The stress-test stage
goes beyond clean-test discrimination by asking whether both ranking and
probability interpretation survive controlled distribution shifts.

At the same time, the implementation is clearly optimized for an
interactive Colab research workflow rather than for a fully packaged
production or benchmark framework. Environment mutation, partially
pinned dependencies, implicit feature-cache validity, and single-seed
secondary analyses limit strict reproducibility. The synthetic
calibration distribution is another major methodological assumption,
particularly because probability calibration can depend on class
prevalence and score coverage.

There are also several nomenclature and documentation issues that should
be corrected before presenting the code as a finalized scientific
benchmark. The \texttt{ece15} naming mismatch should be resolved; the
``multi-scale'' terminology should be stated more precisely; the
packaging verification message should not imply checks that are not
performed; and the reproducibility configuration would benefit from
explicit model and dependency identifiers.

Despite these limitations, the code demonstrates solid research
engineering practices. It does not simply train a model
and report accuracy. Instead, it constructs a reusable feature pipeline,
introduces an explicit leakage boundary, treats calibration as a
first-class problem, evaluates distribution shift, measures confidence
quality with multiple metrics, and preserves intermediate artifacts for
later analysis.

\section{21. Reproducibility Checklist Derived from the
Source}\label{reproducibility-checklist-derived-from-the-source}

A faithful reconstruction of the experiment requires the following
conditions to remain fixed:

\begin{itemize}
\tightlist
\item
  Python and all runtime libraries, especially PyTorch, torchvision,
  NumPy, pandas, scikit-learn, DuckDB, Pillow, and Hugging Face tooling.
\item
  The exact VisA Parquet files identified by their SHA-256 hashes.
\item
  The preprocessing sequence and ImageNet normalization constants.
\item
  The selected backbones and their exact pretrained weight artifacts.
\item
  The reference-bank size \(B=250\), calibration-pool size \(150\), and
  main neighbor count \(k=5\).
\item
  The three primary seeds \(13,42,123\).
\item
  The category configuration containing only \texttt{pcb1}.
\item
  The DINOv2 representation construction from CLS plus mean patch
  tokens.
\item
  L2 normalization before cosine-distance scoring.
\item
  The four synthetic calibration corruptions and their randomization
  rules.
\item
  Logistic regression configuration and isotonic regression settings.
\item
  Adaptive ECE with the implementation's actual default of ten bins.
\item
  The distinction between clean main evaluation, single-seed robustness
  evaluation, and single-seed qualitative diagnostics.
\item
  The feature-cache contents and the path-to-row mapping.
\end{itemize}

For a stronger future reproducibility package, the source should
additionally record explicit environment versions, immutable model
identifiers, data manifests, configuration hashes, and cache provenance.

\section{Conclusion}\label{conclusion}

The provided implementation constitutes a research-oriented
anomaly-detection pipeline built around frozen visual representations, a
normal reference bank, nearest-neighbor anomaly scoring, and post-hoc
probability calibration. The central computational idea is simple but
useful: convert each image into a normalized representation, measure how
close it is to known normal examples, and then learn a separate monotone
mapping from this distance-based score to an anomaly probability.

The implementation adds several layers that make the project more
scientifically interesting than a minimal kNN detector. Dataset
reconstruction is paired with SHA-256 verification, the reference and
calibration pools are separated, synthetic anomalies are introduced
solely for calibration, two distinct calibration models are compared,
calibration quality is evaluated with ECE and Brier score, risk-coverage
behavior is measured, reference size and neighborhood size are varied,
and controlled image corruptions are used to probe distribution-shift
behavior.

The most important interpretive point is that these mechanisms improve
the structure of the experiment, not automatically the strength of its
scientific conclusions. Calibration remains dependent on the synthetic
positive distribution, robustness remains conditional on the selected
corruption family, and reproducibility remains partial because the
environment and model artifacts are not fully pinned. The source is
therefore best understood as a well-structured research prototype whose
strongest contribution is the integration of anomaly scoring,
calibration, and stress testing into one leakage-aware computational
workflow.

For an academic software portfolio, the implementation demonstrates
several valuable capabilities: careful control of data boundaries,
practical use of modern pretrained vision models, numerical feature
processing, non-parametric anomaly scoring, probabilistic calibration,
evaluation design, robustness analysis, artifact management, and
attention to reproducibility. The report shows how those components fit together and,
equally importantly, where the code's methodological assumptions and implementation
constraints should remain visible to a critical technical reader.

\end{document}
